\chapter{IMPLEMENTATION}

This chapter presents the detailed implementation of the proposed intelligent autoscaling framework designed for microservices-based applications. The system integrates real-time monitoring, queuing-theory-based analysis, dynamic scaling, and load balancing into a unified architecture. The implementation focuses on ensuring efficient resource utilization, reduced response time, and stable system performance under dynamic and bursty workloads.

\section{EXPERIMENTAL ENVIRONMENT}

The experimental setup consists of a distributed microservices-based course registration system deployed using Docker Swarm. The system includes multiple independent services such as authentication service, course service, and seat management service. Each service is containerized and deployed as a scalable unit.

An NGINX gateway acts as the entry point for all incoming requests and routes them to the appropriate backend services. Prometheus is integrated into the system for continuous monitoring and collection of real-time metrics. A workload generator is used to simulate real-world traffic conditions, including normal load, peak demand, and sudden burst scenarios. This setup enables comprehensive evaluation of autoscaling performance.\\
\\

\section{DATA COLLECTION AND PROCESSING}

The system continuously collects real-time metrics from all microservices in the distributed architecture using Prometheus as the primary monitoring infrastructure. These comprehensive metrics include fundamental performance indicators such as request arrival rate (requests per second per service), individual service response latency (processing time per request), overall system throughput (successful requests completed per unit time), and active request count (current number of requests being processed). Additional metrics captured include error rates, queue depths and service availability status.

The collected raw data undergoes systematic processing to compute essential queuing-theory parameters that form the mathematical foundation of the intelligent autoscaling system. These parameters include arrival rate ($\lambda$) representing the frequency of incoming requests to each service, service rate ($\mu$) indicating the processing capacity and efficiency of individual service instances, system utilization ($\rho = \lambda/\mu$) showing the current load intensity relative to capacity, and average response delay $D(t)$ calculated using M/M/c queuing formulas. Additional derived metrics include expected queue length ($L_q$) and waiting time in queue ($W_q$), which provide a more accurate and mathematically rigorous representation of system performance compared to traditional resource-based metrics like CPU and memory utilization.

The processed metrics are systematically stored in Prometheus's time-series database and updated at regular intervals (typically every 15-30 seconds) to maintain current system state awareness while ensuring computational efficiency. This continuous data pipeline forms the reliable foundation for all autoscaling decisions, enabling the MFDS algorithm to make informed scaling choices based on actual system behavior rather than reactive threshold violations.

\section{AUTOSCALING MODEL}

Each microservice in the distributed architecture is mathematically modeled as an M/M/c queuing system, representing a multi-server queue with Poisson arrival processes and exponential service times. In this model, $\lambda$ represents the request arrival rate (requests per second) for each service, $\mu$ represents the service rate (requests processed per second per instance), and $c$ represents the current number of active service instances or replicas. This queuing representation captures the fundamental relationship between incoming demand, processing capacity, and system performance in a mathematically rigorous framework.

Using these core parameters, the system calculates critical performance metrics including expected response time using the M/M/c formula $W = W_q + \frac{1}{\mu}$, where $W_q$ is the waiting time in queue, and system utilization $\rho = \frac{\lambda}{c \cdot \mu}$ indicating the load intensity relative to total capacity. The system also computes average queue length $L_q$ and probability of queuing $P(W > 0)$ to provide comprehensive insights into system behavior under current load conditions.

The calculated response delay $D(t)$ is continuously evaluated in real-time and compared against predefined Service Level Agreement (SLA) thresholds defined as the acceptable performance range $[T_{min}, T_{max}]$, where $T_{min}$ represents the lower bound for scale-in decisions and $T_{max}$ represents the upper bound for scale-out actions. Scaling decisions are automatically triggered based on this mathematical comparison: when $D(t) > T_{max}$, the system initiates scale-out operations, and when $D(t) < T_{min}$ with sufficient utilization margin, scale-in operations are considered. This approach ensures that system performance remains within acceptable limits while maintaining optimal resource utilization and cost efficiency.

\section{MICROSERVICE SCALING OPTIMIZATION (MSO)}

\begin{algorithm}[H]
\caption{MSO (Part 1)}
\begin{algorithmic}[1]

\Require $D(t)$, $[T_{min},T_{max}]$, $round\_flag$, $scaling\_flag$, $MS\_Scaling\_List$
\Ensure Updated $MS\_Scaling\_List$

\If{$round\_flag = 0$}
    \For{each microservice $m \in M$}
        \State Calculate scaling factor $R_m$ based on workload
        \State item $\leftarrow (R_m, m)$
        \State Add item to $MS\_Scaling\_List$
    \EndFor
\Else

\If{$scaling\_flag = 1$ and $D(t) > T_{max}$}
    \For{each item in $MS\_Scaling\_List$}
        \State $R_m \leftarrow item[1]$
        \If{$R_m > 0$}
            \State Reduce scaling factor: $R_m \leftarrow \frac{1}{2} R_m$
        \Else
            \State $R_m \leftarrow 0$
        \EndIf
    \EndFor
\EndIf

\end{algorithmic}
\end{algorithm}

\noindent \textit{Algorithm 1 (continued)}

\begin{algorithm}[H]
\caption{MSO (Part 2)}
\begin{algorithmic}[1]

\If{$scaling\_flag = 1$ and $D(t) < T_{min}$}
    \For{each item in $MS\_Scaling\_List$}
        \State $R_m \leftarrow item[1]$
        \If{$R_m > 0$}
            \State Reverse scaling: $R_m \leftarrow -\frac{1}{2} R_m$
        \Else
            \State $R_m \leftarrow 0$
        \EndIf
    \EndFor
\EndIf

\If{$scaling\_flag = -1$ and $D(t) < T_{min}$}
    \For{each item in $MS\_Scaling\_List$}
        \State $R_m \leftarrow item[1]$
        \If{$R_m < 0$}
            \State Adjust scaling: $R_m \leftarrow \frac{1}{2} R_m$
        \Else
            \State $R_m \leftarrow 0$
        \EndIf
    \EndFor
\EndIf

\If{$scaling\_flag = -1$ and $D(t) > T_{max}$}
    \For{each item in $MS\_Scaling\_List$}
        \State $R_m \leftarrow item[1]$
        \If{$R_m < 0$}
            \State Increase scaling: $R_m \leftarrow -\frac{1}{2} R_m$
        \Else
            \State $R_m \leftarrow 0$
        \EndIf
    \EndFor
\EndIf

\State \Return $MS\_Scaling\_List$

\end{algorithmic}
\end{algorithm}
\pagebreak
% =========================
% MFDS (EXPANDED)
% =========================

\section{MULTISTAGE FINE-GRAINED DYNAMIC SCALING (MFDS)}

\begin{algorithm}[H]
\caption{MFDS (Optimized)}
\begin{algorithmic}[1]

\Require Previous deployment $N(t-T)$, delay interval $[T_{min},T_{max}]$
\Ensure Updated deployment $N(t)$

\State Compute response delay $D(t)$

\If{$D(t) \in [T_{min},T_{max}]$}
    \State \Return $N(t-T)$
\EndIf

\If{$D(t) > T_{max}$}
    \State $scaling\_flag \leftarrow 1$
\Else
    \State $scaling\_flag \leftarrow -1$
\EndIf

\State Initialize $MS\_Scaling\_List \leftarrow \emptyset$, $round\_flag \leftarrow 0$

\While{$D(t) \notin [T_{min},T_{max}]$}

    \State $MS\_Scaling\_List \leftarrow$ MSO($\cdot$)
    \State Sort $MS\_Scaling\_List$ by $R_m$
    \State $round\_flag \leftarrow round\_flag + 1$

    \For{each $(R,m)$ in $MS\_Scaling\_List$}
        \While{$R \neq 0$}

            \If{$R < 0$}
                \State Select best host for scale-in
                \State Reduce instances of $m$ and update resources
                \State $R \leftarrow 0$
            \EndIf

            \If{$R > 0$}
                \State Select best host for scale-out
                \State Increase instances of $m$ and update capacity
                \State $R \leftarrow 0$
            \EndIf

        \EndWhile
    \EndFor

    \State Recompute $D(t)$

    \If{$D(t) > T_{max}$}
        \State $scaling\_flag \leftarrow 1$
    \Else
        \State $scaling\_flag \leftarrow -1$
    \EndIf

\EndWhile

\State \Return $N(t)$

\end{algorithmic}
\end{algorithm}

% =========================
% FPR (EXPANDED)
% =========================

\section{FAIR PROBABILISTIC ROUTING (FPR)}

\begin{algorithm}[H]
\caption{FPR (Detailed)}
\begin{algorithmic}[1]

\Require Deployment $N(t)$, network graph $G(V,E)$, request set $F$
\Ensure Routing scheme $R$

\State Initialize routing table $R$
\State Initialize host mapping for all services

\For{each request $f \in F$}

    \State Identify initial host for first service
    \State Store mapping in routing table

    \For{each subsequent service in request chain}

        \State Identify candidate hosts for service
        \State Validate connectivity using graph $G$

        \For{each candidate host}
            \If{no valid connection exists}
                \State Remove host from candidate set
            \EndIf
        \EndFor

        \State Update routing table with valid hosts

    \EndFor

    \State Perform backward validation of routing paths
    \State Remove invalid routing sequences

    \For{each valid host pair}
        \State Compute routing probability based on load and capacity
        \State Assign request probabilistically
        \State Update routing table
    \EndFor

\EndFor

\State \Return routing scheme $R$

\end{algorithmic}
\end{algorithm}

\section{REAL-TIME AUTOSCALING PIPELINE}

The autoscaling pipeline operates as a continuous feedback control system that monitors system metrics, analyzes performance using queuing theory models, and executes intelligent scaling actions in real-time. The pipeline begins with incoming user requests being processed through the NGINX gateway, which serves as the entry point and load balancer for the microservices architecture. As requests flow through the system, Prometheus continuously collects comprehensive real-time metrics from each service including arrival rates, response times, throughput, and queue depths at regular intervals.

The collected metrics feed into the MFDS (Multistage Fine-Grained Dynamic Scaling) algorithm, which serves as the core decision engine that evaluates current system performance against SLA thresholds using M/M/c queuing models. When scaling actions are determined necessary, the MSO (Microservice Scaling Optimization) module computes the optimal number of service replicas required to meet performance objectives, calculating the precise target instance count rather than making incremental adjustments. These scaling decisions are then executed through Docker Swarm's orchestration API, which dynamically adjusts service replica counts across the cluster.

Simultaneously, the FPR (Fair Probabilistic Routing) mechanism operates in parallel to ensure balanced load distribution across all available service instances, dynamically adjusting routing probabilities based on real-time capacity and performance metrics. This integrated pipeline creates a closed-loop control system where monitoring, analysis, decision-making, and execution occur continuously, enabling proactive scaling responses that maintain optimal performance while minimizing resource overhead and ensuring SLA compliance throughout varying workload conditions.

\section{SUMMARY}

This chapter presented the detailed implementation of the intelligent autoscaling system. The integration of real-time monitoring, queuing-theory modeling, and advanced scaling algorithms ensures efficient resource utilization and stable performance. The use of MFDS, MSO, and FPR enables accurate scaling decisions, optimal replica allocation, and balanced traffic distribution under dynamic workloads.